\chapter{RESULTADOS}
\label{cap:resultados}

\section{Experimentos y Resultados}

Los resultados se organizan según los objetivos específicos: \textbf{dos
resultados computacionales} (OE1 y OE2) y una \textbf{salida analítica} (OE3) que
los integra. El punto de partida contra el que se comparan es la línea base
reproducida en la Sección~\ref{sec:linea-base}.

\subsection{Resultado de OE1: modelo híbrido funcional}

Un modelo híbrido verificado de extremo a extremo sobre una sección de control:
la sustitución de $f_{\mathrm{geo}}$ por $f_{\mathrm{ssl}}$ opera dentro del
módulo de prototipos sin romper el razonamiento por hipergrafo y sin degradar el
mIoU \emph{known} respecto de la línea base de la Tabla~\ref{tab:base}. El
entregable incluye la \emph{prueba de funcionamiento} y no solo el código:
inferencia completa, dimensiones documentadas de ambas ramas y verificación de la
compatibilidad de resolución entre el \emph{encoder} PTv3 y la rejilla de
Minkowski.

\subsection{Resultado de OE2: configuración calibrada}

Curvas de sensibilidad de $\alpha$, $M$ y $\epsilon$ sobre la sección de
referencia y una configuración fijada que se aplica sin reajuste al resto de
secciones (Tabla~\ref{tab:sensibilidad}). La lectura esperada es identificar cuál
de los tres parámetros domina el comportamiento del híbrido; en particular, el
valor de $\alpha$ que maximiza el mIoU \emph{novel} indica cuánto peso conserva
la componente geométrica una vez que el descriptor manual ha sido reemplazado.

\begin{table}[H]
\centering
\caption{Sensibilidad de la configuración del híbrido sobre la sección de
referencia.}
\label{tab:sensibilidad}
\begin{tabular}{@{}lccc@{}}
\toprule
\textbf{Parámetro} & \textbf{Rango explorado} & \textbf{Valor fijado} & \textbf{mIoU \emph{novel}} \\
\midrule
$\alpha$   & $[0,\,1]$            & --- & --- \\
$M$        & $\{4,\,8,\,16\}$     & --- & --- \\
$\epsilon$ & \emph{[por definir]} & --- & --- \\
\bottomrule
\end{tabular}
\end{table}

\subsection{Salida analítica de OE3: comparación y ablación}

El resultado central es un incremento medible del mIoU sobre las clases novedosas
respecto de HyperNCD sin modificar, sin degradación de las clases conocidas. La
Tabla~\ref{tab:comparativa} organiza la comparación y la ablación de las tres
variantes del descriptor, aislando la contribución del \emph{encoder}.

\begin{table}[H]
\centering
\caption{Comparación y ablación de las tres variantes del descriptor. Cada fila
reporta además el vector completo de IoU por clase, no solo el agregado
(Sección~\ref{sec:metricas}).}
\label{tab:comparativa}
\begin{tabular}{@{}lcccc@{}}
\toprule
\textbf{Variante del descriptor} & \textbf{mIoU \emph{novel}} & \textbf{mIoU \emph{known}} & $\Delta$\textbf{mIoU} & $G$ \\
\midrule
$[f_{\mathrm{sem}};f_{\mathrm{geo}}]$ (línea base) & --- & --- & --- & 0 \\
$[f_{\mathrm{sem}};f_{\mathrm{ssl}}]$ (propuesto) & --- & --- & --- & --- \\
$[f_{\mathrm{sem}};f_{\mathrm{geo}};f_{\mathrm{ssl}}]$ (Plan~B) & --- & --- & --- & --- \\
\bottomrule
\end{tabular}
\end{table}

\subsection{Entregables}

\begin{enumerate}
\item Línea base reproducida con desviación estándar reportada y análisis de sus
causas, junto con el resultado negativo del Módulo~1
(Sección~\ref{sec:linea-base}).
\item Implementación del modelo híbrido, con prueba de funcionamiento y
repositorio público (OE1).
\item Curvas de sensibilidad de $\alpha$, $M$ y $\epsilon$ y configuración fijada
(OE2).
\item Tablas comparativas, ablación de tres variantes, valores de $G$ y frontera
de operación declarada (OE3).
\end{enumerate}

\section{Discusión}

\paragraph*{Lectura de la métrica.} Un valor de $G$ (Ec.~\eqref{eq:G}) no se
reporta solo. Se acompaña de los tres términos que lo componen y del análisis de
dominancia de la Ec.~\eqref{eq:dominante}: qué clase concreta explica el agregado
y si la ganancia es homogénea entre secciones o está concentrada en una tipología
de elemento. Una mejora que provenga de una sola clase abundante no es el mismo
resultado que una mejora repartida, aunque el mIoU coincida.

\paragraph*{Atribución.} La comparación del \emph{linear probing}
(Ec.~\eqref{eq:probe}) antes de la formación de prototipos con el mIoU final
determina si la ganancia se origina en la representación o en el razonamiento por
hipergrafo. Esa separación es posible porque el \emph{pipeline} es secuencial; si
el sistema fuese unificado, habría que discutir la adecuación misma de la métrica
elegida antes de interpretar el resultado.

\paragraph*{Frontera de operación.} Barriendo densidad de puntos, tamaño de
sección y número de clases novedosas se determina empíricamente la frontera entre
el rango favorable y el desfavorable declarados en el alcance, y se contrasta con
el comportamiento observado en la línea base.

\paragraph*{Realimentación.} Si el análisis muestra que la mejora no se sostiene,
el ajuste se realiza sobre OE1 (esquema de fusión) y OE2 (calibración), no sobre
este capítulo: OE3 es salida analítica y no trabajo computacional independiente.

\paragraph*{Valor para el usuario.} Cada punto de mIoU adicional sobre clases
novedosas es trabajo manual de corrección que deja de realizarse. Los usuarios
están conformes con la precisión que obtienen hoy por vía manual, de modo que la
propuesta de valor no es igualarla sino automatizarla a una precisión suficiente;
sin mejora demostrada sobre el estado del arte, la herramienta no tiene caso.

\paragraph*{Contribución esperada al estado del arte.} La contribución no es una
arquitectura nueva, sino establecer si un hallazgo demostrado en régimen
supervisado ---la nocividad del \emph{geometric shortcut}--- se transfiere al
régimen de descubrimiento de clases novedosas, donde la calidad de la
representación es más crítica porque no hay supervisión sobre las clases
objetivo. Bajo el Plan~B descrito en la Sección~\ref{sec:riesgos}, incluso un
resultado nulo o negativo, documentado con ablación, constituye un hallazgo
reportable.
